\section{RL y Context: por qué la capacidad surge del ciclo cerrado}

Si un modelo solo aprende de un corpus estático, puede absorber los patrones existentes, pero difícilmente sabrá de forma directa si una acción concreta tuvo éxito en una tarea real. El aprendizaje por refuerzo coloca al modelo en un entorno interactivo, donde intenta acciones, observa los resultados y actualiza su política según la reward. Aquí, \term{RL} es Reinforcement Learning, que en español suele traducirse como aprendizaje por refuerzo; \term{environment} es el mundo de la tarea en el que el modelo puede actuar y recibir retroalimentación sobre el estado o el resultado. Ambos amplían «más tokens» a «más experiencia verificada».

\begin{figure}[H]
\centering
\includegraphics[width=0.9\textwidth]{images/00-19-45-how-rl-works.jpg}
\caption{Ciclo cerrado del aprendizaje por refuerzo: el entorno entrega un estado, la política ejecuta una acción y la retroalimentación impulsa la actualización de la política.}
\figsource{Video oficial de Stanford Online 00:19:45, `images/00-19-45-how-rl-works.jpg`.}
\end{figure}

\begin{knowledgebox}{Lectura de la figura: separar los tres tipos de objetos}
En la figura, el agent o policy es el modelo que toma decisiones, el environment es el sistema externo que responde a las acciones y la reward es una retroalimentación comprimida sobre el resultado. Ninguno puede sustituir a los otros: un modelo más grande no equivale a un mejor entorno, un entorno más rico no garantiza que la reward sea correcta, y que la reward pueda calcularse no significa que represente el resultado que las personas realmente quieren. La capacidad del sistema surge de la calidad del ciclo cerrado, no de la escala aislada de alguno de sus componentes.
\end{knowledgebox}

Un objetivo de RL simplificado puede escribirse como
\[
J(\theta)=\mathbb{E}_{\tau\sim\pi_\theta}\left[\sum_{t=0}^{T}\gamma^t r_t\right].
\]
Aquí, $\pi_\theta$ denota la política con parámetros $\theta$, $\tau$ es una trajectory que surge de la interacción entre la política y el entorno, $r_t$ es la reward del paso $t$, $\gamma$ es el factor de descuento de los rendimientos futuros y $T$ es el instante en que termina la tarea. La fórmula no responde cómo debe diseñarse la reward ni garantiza que el entrenamiento sea estable; solo expresa como objeto de optimización «hacer más probables las trayectorias de alto rendimiento».

\begin{warningbox}{Nota de clase: el scaling empírico no es scaling ilimitado}
El ponente describe una y otra vez las scaling laws como empirical observations. Que el desempeño haya mejorado de forma sostenida al aumentar la escala de entrenamiento solo indica que existe una regularidad dentro del intervalo observado. El RL también puede estancarse por la calidad del entorno, fallas en el verificador, exploración insuficiente, cambios de distribución u objetivos mal alineados. La planeación de ingeniería puede extrapolar curvas, pero debe preparar al mismo tiempo las condiciones de falla y experimentos alternativos.
\end{warningbox}

\subsection{Basic AI Scaling Recipe}

Esta sección reduce el proceso de producción de las capacidades de los modelos modernos a cuatro nodos: el compute sostiene la research, la research produce models, los models llegan a los usuarios y las tareas mediante la inference, y el uso genera a su vez context y retroalimentación. Lo importante de este diagrama no es el número de flechas, sino la dirección del ciclo cerrado. El entrenamiento no encapsula el conocimiento en los pesos de una sola vez; el entorno después del despliegue, las llamadas a herramientas y los registros de éxitos y fracasos determinan lo que la siguiente ronda de investigación y entrenamiento podrá ver.

\begin{figure}[H]
\centering
\includegraphics[width=0.76\textwidth]{images/00-22-34-basic-ai-scaling-recipe.jpg}
\caption{Basic AI Scaling Recipe: el ciclo de compute, research, models, inference y context feedback.}
\figsource{Video oficial de Stanford Online 00:22:34, `images/00-22-34-basic-ai-scaling-recipe.jpg`.}
\end{figure}

\begin{importantbox}{Lectura de la figura: la flecha del ciclo que más se pasa por alto}
La ruta de compute a model es fácil de cuantificar: el número de GPU, los FLOPs de entrenamiento y el tamaño del modelo tienen cifras. La ruta de regreso de inference a context feedback es más difícil: exige saber en qué flujo de trabajo usan los usuarios el modelo, cuándo falla, si la falla puede verificarse, si hay autorización para que los datos regresen y si la retroalimentación representa a la población objetivo. El ponente coloca el context en primer lugar porque esta flecha suele determinar si el compute posterior tendrá usos de alto valor.
\end{importantbox}

\subsection{El context no es solo un prompt más largo}

Aquí, el alcance del context es mucho mayor que la ventana de conversación. Incluye la descripción de la tarea, los documentos privados, el código fuente, los permisos de herramientas, el entorno de ejecución, las decisiones históricas, el comportamiento de los usuarios y los resultados verificables. Aun con los mismos pesos, un modelo que difiere en si puede acceder a la terminal, a las pruebas, a los registros de despliegue y al conocimiento de la organización se comporta como un sistema completamente distinto. Por eso Frontier Systems considera el context un activo de infraestructura: hay que recopilarlo, autorizarlo, depurarlo, conectarlo y mantenerlo de forma continua.

\begin{figure}[H]
\centering
\includegraphics[width=0.9\textwidth]{images/00-24-04-context-environments-flywheel.jpg}
\caption{Context/Environments flywheel: integrarse al flujo de trabajo, observar las fallas, mejorar el entorno y el modelo, y ampliar el uso.}
\figsource{Video oficial de Stanford Online 00:24:04, `images/00-24-04-context-environments-flywheel.jpg`.}
\end{figure}

\begin{knowledgebox}{Lectura de la figura: tres preguntas determinan si el flywheel funciona}
Primero, si el sistema entra en contacto con tareas reales de alto valor o solo gira sobre datos de demostración; segundo, si el éxito y el fracaso pueden verificarse de forma confiable o solo se cuenta con una satisfacción difusa; tercero, si el regreso de los datos cuenta con la autorización del usuario y respeta los límites de seguridad. El flywheel funciona solo cuando el crecimiento del uso, la calidad de la retroalimentación y la mejora del modelo se refuerzan entre sí. Aumentar únicamente el volumen de llamadas puede amplificar el ruido, los riesgos de privacidad y la automatización de errores.
\end{knowledgebox}

\begin{table}[H]
\centering
\begin{tabular}{L{0.18\textwidth}L{0.32\textwidth}L{0.4\textwidth}}
\toprule
Término & Problema que resuelve & Significado en esta clase \\
\midrule
Prompt context & Qué información necesita la solicitud actual & Conversación, documentos y restricciones inmediatas \\
Tool context & Qué acciones puede realizar el modelo & API, terminal, búsqueda, bases de datos y permisos \\
Environment & Dónde se ejecuta la acción y produce resultados & Repositorio de código, navegador, laboratorio o sistema de negocio \\
Feedback & Cómo juzgar el efecto de una acción & Pruebas, modificaciones del usuario, ingresos, latencia o incidentes de seguridad \\
Flywheel & Por qué las mejoras pueden acumularse de forma continua & El uso genera retroalimentación de alta calidad, y esta mejora el modelo y el producto \\
\bottomrule
\end{tabular}
\caption{Glosario: descomponer el context genérico en componentes que pueden diseñarse y auditarse.}
\end{table}

\subsection{Resumen de la sección}

El RL ofrece un ciclo cerrado de aprendizaje política–entorno–retroalimentación, y la Basic Scaling Recipe sitúa ese ciclo dentro del proceso full-stack de compute, research, model e inference. Lo decisivo del context no es su longitud, sino la calidad de las tareas, las herramientas, los entornos y las señales de verificación. Y precisamente porque el context de alto valor es escaso, las empresas de modelos y las empresas de aplicaciones compiten intensamente por los puntos de acceso, los datos y los flujos de trabajo.
